\section{Experimental Setup}
\label{sec:setup}

\subsection{Single harness, reproducible reference}

All results come from one implementation, which reproduces the reference
network exactly. We verified this by running the reference configuration
(\texttt{--activation poly\_proj}, i.e.\ Eq.~\eqref{eq:projective-poly}) for
two epochs and confirming that the per-epoch train/test accuracies agree with
the standalone reference script to the last reported digit. The fidelity run's
values are $\fidEOneTrain\%$ / $\fidEOneTest\%$ after epoch one and
$\fidETwoTrain\%$ / $\fidETwoTest\%$ after epoch two (train / test); note that
the reference script's default epoch count is twenty, and the cosine schedule
has period equal to the epoch budget, so the two-epoch fidelity check is
compared against the reference script run with the same two-epoch budget.
Because the harness is a strict superset of the reference script, every number
in this paper is comparable to every other number, and the reference
configuration appears in every table as the baseline row: with its own
twenty-epoch budget it reaches $\refTrainAcc\%$ train and $\refTestAcc\%$ test
accuracy in $\refEpochSec$\,s per epoch ($\refTotalSec$\,s total).

\subsection{Data}

\paragraph{MNIST.} $60{,}000$ training and $10{,}000$ test grayscale images,
$28 \times 28$, flattened to $784$ dimensions with the raw pixel scale in
$[0,1]$ (only \texttt{ToTensor}, no mean-variance standardization --- the
reference pipeline uses none, and adding one would change the effective scale
of the activation, which is precisely one of the axes under study). Ten
classes.

\paragraph{Synthetic.} Eight generator families, each a deterministic function
of a seed, with the test split drawn from the same generator at a different
seed so that the two splits are independent samples of the same distribution
rather than a split of one finite sample.

\begin{itemize}
\item \emph{rings} (5 classes, 2-D): concentric noisy circles; requires the
  network to use the radius, i.e.\ a \emph{quadratic} function of the input.
  Directly relevant to whether a cubic activation helps on a polynomial task.
\item \emph{spiral} (5 classes, 2-D): the classical interleaved spiral;
  requires a strongly nonlinear, non-monotone decision boundary.
\item \emph{blobs} (5 classes, 2-D): well-separated Gaussians; an easy control
  that should saturate for every activation and therefore calibrates the other
  results.
\item \emph{xor} (2 classes, 2-D): parity of the coordinate signs; the minimal
  case where an additive (linear-plus-pointwise) model must build a product.
\item \emph{checker} (2 classes, 2-D): a $4\times4$ checkerboard; high-frequency
  and piecewise constant, the regime where a smooth activation is expected to
  have difficulty.
\item \emph{cubes} (2 classes, 3-D): parity of three dichotomized coordinates,
  i.e.\ a \emph{third-order} interaction, matched to the order of the reference
  cubic.
\item \emph{parity} (2 classes, 10-D): parity of ten binary features. The label
  depends on the tenth-order interaction, so any model that can only form
  low-order or additive functions must fail; this is our stress test for
  whether a degree-3 activation's advantage is about order or about scale.
\item \emph{gauss\_highdim} (5 classes, 50-D): a five-class problem living on a
  two-dimensional subspace, embedded in fifty dimensions with additive noise.
  Tests whether the projection's scale control helps when most input
  directions are noise.
\end{itemize}

For the synthetic experiments the input is standardized using \emph{training}
split statistics only. We use $4000$ training and $800$ test samples per
family.

\subsection{Models and hyperparameters}

Unless stated otherwise, the configuration is depth $L = 10$ linear layers,
hidden width $H = 128$, AdamW with learning rate $10^{-3}$ and weight decay
$10^{-4}$, cosine-annealed over the run, gradient-norm clipping at $5.0$,
batch size $128$, and $10$ epochs. The learning rate, weight decay, clipping
threshold, batch size and scheduler all match the reference script, so the
reference row of each table is a faithful baseline rather than a re-tuned one.
The synthetic experiments use the shallower, smaller setting requested for
that part of the study: $L = 3$, $H = 64$, learning rate $3\times 10^{-3}$,
batch size $64$, $60$ epochs. The headline comparison is repeated over five
seeds ($0$--$4$) and reported as mean $\pm$ standard deviation; all other
tables use the single reference seed $42$.

\subsection{Sweeps}

\begin{table}[t]
\centering
\caption{Sweeps. Each row varies one factor of the reference configuration
while holding the rest at the reference values. $L$ = depth, $H$ = width,
$d$ = polynomial order, $\lambda$ = polynomial coefficient, pos = placement
of the normalization relative to the nonlinearity.}
\label{tab:sweeps}
\small
\begin{tabular}{@{}llll@{}}
\toprule
Axis & Values & Fixed & Purpose \\
\midrule
Activation (pointwise) & 6 functions & $L{=}10$, no norm & main comparison \\
Activation (gated) & 3 units, $W_g{=}82$ & $L{=}10$, no norm & modern LLM shapes \\
Learned polynomial & Legendre/Chebyshev/spline & $L{=}10$ & fixed vs.\ learned \\
Polynomial order & $d \in \{2,3,4,5\}$ & $\lambda{=}10^{-2}$, $\pi$ & order sensitivity \\
Coefficient & $\lambda \in \{10^{-3},10^{-2},5{\cdot}10^{-2},10^{-1}\}$ & $d{=}3$, $\pi$ & coefficient sensitivity \\
Normalization & none, $\pi$, BN, LN, RMS & ReLU \& poly & main comparison \\
Placement & pre, post & ReLU \& poly & pre- vs.\ post-norm \\
Depth & $L \in \{2,3,5,10,15,20\}$ & 7 activation/norm pairs & depth scaling \\
Depth (extreme) & $L \in \{30,50\}$ & 7 activation/norm pairs & gradient propagation \\
Width & $H \in \{82,128,192\}$ & 4 activations & capacity control \\
\bottomrule
\end{tabular}
\end{table}

Table~\ref{tab:sweeps} lists the axes. Three design choices deserve comment.
First, wherever an architectural family changes the parameter count --- gated
units, learnable activations, the spline basis --- we report the matched
parameter count alongside the accuracy, and we include a width-controlled
pointwise baseline, so that no conclusion rests on an accidental capacity
difference. Second, following the same deployability criterion that fixed the
activation list, we restrict the sweep to activations that a current frontier
language model actually uses, plus the polynomial family that is the subject
of the paper; the historical activation zoo is not re-audited. Third, the
depth sweep includes $30$ and $50$ layers, well past the depth at which
unnormalized networks are traditionally expected to fail, because the claim
most worth testing about a scale-controlling activation is whether it keeps
very deep stacks alive.

\subsection{Diagnostics and cost accounting}

For every run we record per-layer gradient norms and per-layer activation
norms before training and after the last epoch, on a fixed mini-batch of $512$
training samples. We also record parameter count, forward FLOPs per sample,
and wall-clock seconds per epoch. All runs execute on CPU (the reference
environment has no GPU), with the thread count pinned and configurations run
in parallel processes; wall-clock numbers are therefore comparable within the
study but are not a statement about GPU throughput. Every configuration writes
a JSON record containing its arguments, metrics, per-epoch history, per-class
accuracy and confusion matrix, so all tables and figures in this paper are
generated mechanically from those records by a separate analysis script
rather than transcribed by hand.
